% Exported from Math++. Compile with pdflatex.
\documentclass[11pt,a4paper]{article}
\usepackage[T1]{fontenc}
\usepackage[utf8]{inputenc}
\usepackage{lmodern}
\usepackage{amsmath,amssymb,amsthm,mathtools}
\usepackage{graphicx,booktabs,xcolor}
\usepackage[hidelinks]{hyperref}
\usepackage[margin=25mm]{geometry}
\theoremstyle{plain}
\newtheorem{theorem}{Theorem}
\newtheorem{lemma}[theorem]{Lemma}
\newtheorem{corollary}[theorem]{Corollary}
\newtheorem{proposition}[theorem]{Proposition}
\newtheorem{conjecture}[theorem]{Conjecture}
\theoremstyle{definition}
\newtheorem{definition}[theorem]{Definition}
\newtheorem{example}[theorem]{Example}
\theoremstyle{remark}
\newtheorem*{remark}{Remark}
\setlength{\parindent}{0pt}\setlength{\parskip}{0.6em}
\begin{document}
\title{Proofs, step by step}
\author{Math++ examples}
\date{}
\maketitle
Each step of a proof is a claim and a reason. The checker marks every step in the equation view (checked symbolically, checked numerically, checked by Lean, structure only, cited, a gap or an external claim, or wrong) and gives each theorem a summary that never claims more than its weakest step. None of it is written into this file.

\begin{lemma}
If $p$ is odd, then $p^2$ is odd.

\end{lemma}
\begin{proof}
\begin{enumerate}
\item Write $p = 2m + 1$ for an integer $m$ --- by assumption.
\item $p^2 = 4m^2 + 4m + 1 = 2(2m^2 + 2m) + 1$ --- by algebra.
\item So $p^2$ is odd, and the assumption is used up --- by cases.
$\blacksquare$
\end{enumerate}
\end{proof}
\begin{theorem}
For every integer $n \ge 1$, $\sum_{k=1}^{n} k = \frac{n(n+1)}{2}$.

\end{theorem}
\begin{proof}
\begin{enumerate}
\item Base case: $\sum_{k=1}^{1} k = 1 = \frac{1 \cdot 2}{2}$ --- by algebra.
\item Inductive step: if $\sum_{k=1}^{n} k = \frac{n(n+1)}{2}$, then $\sum_{k=1}^{n+1} k = \frac{n(n+1)}{2} + (n+1) = \frac{(n+1)(n+2)}{2}$ --- by algebra.
\item So the formula holds for every $n \ge 1$ --- by induction.
$\blacksquare$
\end{enumerate}
\end{proof}
\begin{theorem}
For every real $x$, $(x + 1)^2 = x^2 + 2x + 1$.

\end{theorem}
\begin{proof}
\begin{enumerate}
\item For every real $x$: $(x + 1)^2 = x^2 + 1$ --- by algebra.
$\blacksquare$
\end{enumerate}
\end{proof}
\begin{theorem}
Every integer $N > 1$ has a prime factor.

\end{theorem}
\begin{proof}
\begin{enumerate}
\item This is Elements VII.31 --- \emph{(cited: every integer greater than 1 has a prime factor \cite{euclid})}.
$\blacksquare$
\end{enumerate}
\end{proof}
\begin{theorem}
The tail vanishes.

\end{theorem}
\begin{proof}
\begin{enumerate}
\item The boundary term is zero --- by the divergence theorem \emph{(gap: decay at infinity is assumed)}.
$\blacksquare$
\end{enumerate}
\end{proof}
\begin{conjecture}
Every even number above 2 is a sum of two primes.

\end{conjecture}
\subsection*{References}

\begin{itemize}
\item \cite{euclid} Euclid, \emph{Elements}, Book VII Prop. 31.
\end{itemize}
\end{document}
